\chapter{Temu、SHEIN、TikTok Shop 平台模式}
\label{ch:platform}

\begin{insight}[本章要解决的问题]
\begin{enumerate}
  \item 同一个工厂、同一条产线，把货供给 Temu 全托管、供给 SHEIN 的第三方店、
        还是拿到 TikTok Shop 自己开播，为什么净利润能差出两到三倍？
        差的到底是售价，还是\kw{风险与资金的归属}？
  \item 「全托管」听起来是把定价、物流、售后全部外包，卖家只负责生产——
        那卖家到底还承担了什么风险？为什么很多工厂做着做着，
        发现自己成了平台的\kw{无息资金方}与\kw{质量赔付的承担人}？
  \item 平台的\kw{核价}到底在比什么？是比同款最低价，还是比「同一价格带里最能打的那一个」？
        卖家报价的可行域边界在哪里，为什么「成本加成定价」在这个体系里几乎必然亏钱？
  \item SHEIN 的「小单快反」为什么同行学了十年仍然学不像？
        它依赖的是软件系统、产业带密度，还是买手组织与账期结构？
        首单到底该下多少件，返单的触发条件怎么算？
  \item TikTok Shop 的「内容--交易一体」在美区与东南亚为什么是两套完全不同的打法？
        达人分销的佣金率定到多少才是盈亏平衡？直播间的「人货场」拆开来对应哪些可优化变量？
  \item 平台的罚款、下架、冻结货款与封店，从规则条文到实际执行之间有多大的解释空间？
        卖家有没有可操作、可复盘的申诉与资金解冻路径？
  \item 面对「规则随时会变」的平台，卖家应该做加法（多平台、多模式并行）
        还是做减法（单点做到极致）？两者的财务含义与风险结构有什么不同？
\end{enumerate}
\end{insight}

Temu、SHEIN、TikTok Shop 常被放在一起讨论，但它们做的其实不是同一件事。
如果只用一个视角来看，最有解释力的说法是：
\kw{这三家都在把中国供应链的产能，用算法与跨境物流直接接到海外消费者}，
把过去「工厂--外贸公司--进口商--批发商--零售商」这条链上被层层加价的环节，
压缩成「工厂--平台--消费者」的一次性加价。
\cite{statista_ecommerce} 一类公开市场数据反复显示，跨境电商的增长主要来自中间环节的压缩，
而不是终端售价的抬升；消费者付得更少，工厂拿到的单价也未必更高，
差额大部分被\kw{履约成本与流量成本}吃掉了。

但三家压缩的方式完全不同。
Temu 走的是\kw{平台接管需求侧}：平台定价、平台投流、平台履约，卖家退化成产能供给方，
核心竞争维度只剩「报价」和「交期」。
SHEIN 走的是\kw{平台接管供给侧}：用自有的柔性供应链体系把「预测式大批量生产」
换成「小批量试单加数据驱动返单」，用极短的库存周转天数换取资金效率。
TikTok Shop 走的是\kw{把需求创造与成交合并在同一个内容场里}：
短视频与直播既是广告也是货架，达人既是渠道也是流量合伙人。

这个区别决定了三件事，也构成本章的主线：
第一，\kw{责任边界在哪里}——谁定价、谁履约、谁承担退货与库存；
第二，\kw{利润从哪里来}——是供货价差、是资金周转，还是内容转化；
第三，\kw{风险由谁承担}——被压价、被罚款、被冻结货款，还是内容衰退与账号风险。

技术性约定：本章涉及的所有平台规则、费率、佣金比例、罚款标准、账期长度、
入驻门槛与处罚机制，都属于\kw{变动极快}的政策性内容，
一律\textbf{以各平台官方最新规则与卖家后台公告为准}；
本章给出的数值只作为「量级参考」与「建模演示」，不构成任何合规或投资建议。
凡是引用公开报道的数据，都会明确标注为「公开报道口径，存在不确定性」；
本章不编造任何未公开的内部经营数据或财务数字。
涉及跨境经营的税务、数据与广告合规问题，读者必须以目的国法律与专业机构意见为准，
相关全景可回看 \cref{ch:cbec}。

\section{三种模式：全托管、半托管与自主运营}
\label{sec:plat-modes}

\subsection{模式的定义：决策权在谁手里}

跨境电商平台与卖家之间的合作模式，可以用一个非常朴素的问题切开：
\kw{从商品离开工厂到消费者收到货，中间的每一个决策由谁做、每一个损失由谁赔}。
按这个标准，行业里通行的三种模式定义如下（具体命名与细则各平台不同，以官方文档为准）。

\begin{definition}[全托管模式]
平台（或其关联采购主体）向卖家\emph{采购}商品，卖家只负责按约定标准生产并把货送到平台指定的境内集货仓；
平台负责\kw{定价、营销投流、跨境干线、海外仓、尾程配送、售后与退款}。
卖家是供货商（supplier），收入形式是\kw{供货价}结算，而不是零售价分成。
\end{definition}

\begin{definition}[半托管模式]
卖家自己负责\kw{海外仓库存与本地履约}（备货到海外仓并完成本地发货），
平台仍然掌握\kw{定价权、前台流量分配与大部分营销投放}，并对卖家报价进行核价；
平台负责跨境干线（或由卖家自行承担头程，视平台细则而定）与消费者侧的交易体验。
卖家是「半个商家、半个供货商」。
\end{definition}

\begin{definition}[自主运营模式]
卖家用自有主体在平台开店（或运营独立站与内容账号），
\kw{自主定价、自主投流、自主履约、自主售后}，
平台只提供交易场所、支付、流量工具与规则框架，并从成交额中抽取佣金。
卖家是真正的商家（merchant）。
\end{definition}

这三种定义里最关键的一句话是：\textbf{全托管卖的是「产能」，半托管卖的是「海外库存运营能力」，自主运营卖的是「需求获取能力」}。
三者对卖家的能力要求不是程度差异，而是\kw{种类差异}——
一个擅长接单排产的工厂，未必有能力运营海外仓，更未必有能力在内容平台上获取流量。

\begin{table}[htbp]
\centering
\caption{三种模式的责任划分：定价权、物流、售后与流量}
\label{tab:plat-responsibility}
\small
\begin{tabularx}{\textwidth}{L{2.1cm} L{3.5cm} L{3.5cm} Y}
\toprule
\tabhead{责任维度} & \tabhead{全托管} & \tabhead{半托管} & \tabhead{自主运营} \\
\midrule
定价权 & 平台核价并最终定价，卖家只报供货价 & 平台定价并核价，卖家报价空间略大 & 卖家自主定价 \\
\addlinespace
国内段物流 & 卖家送到平台国内集货仓 & 卖家自行送仓或自建头程 & 卖家自行组织 \\
\addlinespace
跨境干线与清关 & 平台 & 平台或卖家（看细则） & 卖家（或委托货代） \\
\addlinespace
海外仓与尾程 & 平台 & 卖家（本地发货能力是门槛） & 卖家，或用平台仓配服务 \\
\addlinespace
售后与退款 & 平台对消费者兜底，再按规则向卖家追偿 & 平台规则下由卖家承担主要部分 & 卖家全额承担 \\
\addlinespace
流量获取 & 平台 100\% 负责，卖家不投流 & 平台为主，部分平台允许卖家自投 & 卖家自投广告与内容 \\
\addlinespace
库存风险 & 卖家承担（货已交付但可能滞销、滞销后退回） & 卖家承担（海外仓滞销成本更高） & 卖家承担 \\
\addlinespace
卖家角色 & 供货商 & 半供货商半商家 & 商家 \\
\addlinespace
结算对象 & 供货价（平台采购价） & 结算价（扣佣后的成交额） & 结算价（扣佣、扣支付费后） \\
\bottomrule
\end{tabularx}
\end{table}

\begin{figure}[htbp]
\centering
\begin{tikzpicture}[
  every node/.style={font=\sffamily\tiny, align=left},
  colbox/.style={draw=NavyMid, fill=BoxBG, rounded corners=2pt, inner sep=5pt,
                 text width=4.1cm},
  colboxb/.style={draw=TealDeep, fill=BoxBGAlt, rounded corners=2pt, inner sep=5pt,
                  text width=4.1cm},
  colboxc/.style={draw=AmberDeep, fill=white, rounded corners=2pt, inner sep=5pt,
                  text width=4.1cm},
  tit/.style={font=\sffamily\small\bfseries\color{NavyDeep}},
  note/.style={draw=RuleGray, fill=BoxBGAlt, rounded corners=2pt, inner sep=4pt,
               text width=13.6cm, align=left, font=\sffamily\tiny}
]
\node[tit] at (0,0.55) {全托管};
\node[tit] at (4.8,0.55) {半托管};
\node[tit] at (9.6,0.55) {自主运营};
\node[colbox] (a) at (0,-2.5)
  {\textbf{定价}：\textcolor{NavyMid}{平台}\\[2pt]
   \textbf{海外物流}：\textcolor{NavyMid}{平台}\\[2pt]
   \textbf{售后}：\textcolor{NavyMid}{平台兜底 + 向卖家追偿}\\[2pt]
   \textbf{流量}：\textcolor{NavyMid}{平台 100\%}\\[2pt]
   \textbf{库存风险}：\textcolor{AmberDeep}{卖家}\\[2pt]
   \textbf{核心能力}：\textcolor{AmberDeep}{成本与交期}};
\node[colboxb] (b) at (4.8,-2.5)
  {\textbf{定价}：\textcolor{NavyMid}{平台}（核价）\\[2pt]
   \textbf{海外物流}：\textcolor{AmberDeep}{卖家本地发货}\\[2pt]
   \textbf{售后}：\textcolor{AmberDeep}{卖家为主}\\[2pt]
   \textbf{流量}：\textcolor{NavyMid}{平台为主}\\[2pt]
   \textbf{库存风险}：\textcolor{AmberDeep}{卖家（海外仓）}\\[2pt]
   \textbf{核心能力}：\textcolor{AmberDeep}{海外仓运营}};
\node[colboxc] (c) at (9.6,-2.5)
  {\textbf{定价}：\textcolor{AmberDeep}{卖家}\\[2pt]
   \textbf{海外物流}：\textcolor{AmberDeep}{卖家}\\[2pt]
   \textbf{售后}：\textcolor{AmberDeep}{卖家全额}\\[2pt]
   \textbf{流量}：\textcolor{AmberDeep}{卖家自投 + 内容}\\[2pt]
   \textbf{库存风险}：\textcolor{AmberDeep}{卖家}\\[2pt]
   \textbf{核心能力}：\textcolor{AmberDeep}{需求获取}};
\node[note] at (4.8,-5.6)
  {从全托管到自主运营，卖家的\textbf{决策权}上升、\textbf{毛利上限}上升，
   但\textbf{前置资金}、\textbf{专业能力门槛}与\textbf{经营风险}同时上升。
   选型的唯一正确起点是：先盘清自己已经具备的能力，再倒推模式，而不是反过来。};
\end{tikzpicture}
\caption{全托管、半托管与自主运营的责任边界对比}
\label{fig:plat-responsibility}
\end{figure}

\subsection{卖家角色与结算方式：供货价、结算价与账期}

模式差异最终会落到\kw{结算口径}上，而这是新卖家最容易搞混的地方。

\term{供货价}{Supply Price} 是全托管模式下平台的采购价，
它在形式上与外贸订单的 FOB 价（贸易术语见 \cite{incoterms2020}）很像，但有两个本质区别：
第一，它通常由平台\kw{核价}产生，而不是买卖双方议价；
第二，它一般\kw{不含任何营销费用与履约费用}，但这些费用的风险却通过罚款、
退货追偿、滞销退回等方式部分回传到卖家身上。

\term{结算价}{Settlement Price} 是自主运营（与半托管）模式下的口径：
消费者支付额扣除平台佣金、支付通道费、退款与折扣补贴之后，
真正打到卖家账户的金额。它看起来比供货价高很多，
但卖家要用它去覆盖头程、关税、海外仓、尾程、广告、退货与人工。

把单件利润写成可计算的形式，全托管模式下近似为
\begin{equation}
\pi_{\text{全托管}}
= p_{w} - c_{\text{物料}} - c_{\text{加工}} - c_{\text{包装}}
- c_{\text{国内干线}} - \lambda_{\text{赔付}}\, p_{w}
- \kappa_{\text{退货}}\, c_{\text{退货}} - \frac{F_{\text{固定}}}{Q} ,
\label{eq:fm-profit}
\end{equation}
其中 $p_{w}$ 是供货价，$\lambda_{\text{赔付}} \in [0,1]$ 是因品控、超时、缺货、
虚假宣传等原因产生的\kw{赔付与罚款占供货额的比例}，
$\kappa_{\text{退货}}$ 是退货率（全托管下退货运费多由平台承担，
但货值损失与二次销售折损常由卖家吸收），
$F_{\text{固定}}/Q$ 是打样、认证、模具、质检等固定费用在批量 $Q$ 上的摊销。
自主运营模式下的同款商品则近似为
\begin{equation}
\pi_{\text{自运营}}
= p_{r}\,(1 - r_{\text{平台}} - \rho_{\text{退款}})
- c_{\text{货}} - c_{\text{头程}} - c_{\text{关税}} - c_{\text{仓配}}
- c_{\text{获客}} - c_{\text{售后}} ,
\label{eq:fm-profit-self}
\end{equation}
其中 $p_{r}$ 是零售价，$r_{\text{平台}}$ 是平台综合佣金与支付费，
$\rho_{\text{退款}}$ 是退款率，$c_{\text{获客}}$ 是摊到每单的获客成本
（广告费除以订单数，即 $\CAC$ 的单量口径，详见 \cref{ch:ad}）。
\cref{eq:fm-profit,eq:fm-profit-self} 放在一起看，就能解释一个反复出现的现象：
\textbf{全托管的毛利结构像代工，自主运营的毛利结构像品牌零售}——
前者的天花板由成本与报价决定，后者的天花板由转化率与复购决定。

\begin{keypoint}[账期才是全托管最容易忽视的成本项]
全托管模式下，卖家通常需要\kw{先生产、先交货、后结算}，
平台按约定周期（常见说法是「发货后若干天」或「签收后若干天」，
具体以各平台官方结算规则为准）支付货款。
这意味着卖家的\kw{现金转换周期}（Cash Conversion Cycle）被显著拉长。
用经典的营运资金恒等式表达：
\begin{equation}
\text{现金缺口} \;\approx\; \frac{\text{年营业成本}}{365}
\cdot \underbrace{\big( D_{\text{库存}} + D_{\text{应收}} - D_{\text{应付}} \big)}_{\text{现金转换周期}},
\label{eq:fm-cashgap}
\end{equation}
其中 $D_{\text{库存}}$ 是原材料与成品平均在库天数，
$D_{\text{应收}}$ 是从交货到收款的账期，$D_{\text{应付}}$ 是对上游供应商的平均付款账期。
\textbf{当一个卖家对上游是现款现货（$D_{\text{应付}} \approx 0$），
而平台账期是 30 到 45 天（$D_{\text{应收}} = 30$ 到 $45$）时，
他的现金缺口几乎等于整条供应链的资金需求}。
把这个缺口折成年化资金成本加进去，
「看起来有 15\% 毛利」的全托管生意，可能只剩下个位数的净利。
这就是为什么同一款货，\kw{自有资金充裕的工厂能接，靠借贷周转的贸易商接不了}。
\end{keypoint}

\subsection{三种模式对卖家能力的要求与利润结构}

\begin{table}[htbp]
\centering
\caption{三种模式的能力要求、利润结构与风险暴露（利润区间为示范量级，非真实数据）}
\label{tab:plat-capability}
\small
\begin{tabularx}{\textwidth}{L{2.0cm} L{3.4cm} L{3.4cm} Y}
\toprule
\tabhead{对比项} & \tabhead{全托管} & \tabhead{半托管} & \tabhead{自主运营} \\
\midrule
核心能力 & 成本控制、交期稳定、品控一致 & 海外仓库存管理与补货预测 & 选品、内容、投放、复购运营 \\
\addlinespace
最不需要的能力 & 流量运营、海外客服 & 前台流量运营 & 无，全部需要 \\
\addlinespace
定价权 & 无（只能报价） & 很弱 & 完整 \\
\addlinespace
毛利结构 & 供货价减成本，毛利率常见为个位数到二十几个百分点（示范量级） & 结算价减货值与仓配，毛利率常高于全托管 & 零售价减全链路成本与获客成本，毛利率最高但波动最大 \\
\addlinespace
资金占用 & 中（备货到国内仓） & 高（海外仓压货） & 高（备货加广告预付） \\
\addlinespace
主要风险 & 被压价、罚款、滞销退回 & 海外滞销、补货判断错误 & 内容衰退、广告成本上升、账号与合规风险 \\
\addlinespace
起量速度 & 快（平台推流） & 中 & 慢，取决于内容与投放能力 \\
\addlinespace
可累积资产 & 几乎没有（无用户、无品牌） & 少量（本地履约能力） & 用户、品牌、私域、达人关系 \\
\bottomrule
\end{tabularx}
\end{table}

\begin{insight}[「利润结构」比「利润率」更值得比较]
多数卖家比较模式时只看利润率，这是错的。
更完整的比较应该看\kw{利润的三种来源}：
(1) \textbf{生产性利润}——来自成本控制与规模效应，全托管主要靠这个，
但它有硬天花板，因为成本下降空间是有限的；
(2) \textbf{周转性利润}——来自资金周转速度，
在单件利润率 5\% 但一年周转 12 次的情形下，资金回报率仍然可观
（这正是 SHEIN 逻辑的核心，见 \cref{eq:shein-turn}）；
(3) \textbf{需求性利润}——来自流量获取效率与品牌溢价，自主运营靠这个，
它上限最高，但也最容易被平台规则与算法变化摧毁。
一个健康的卖家应该同时拥有至少两种利润来源，
只依赖一种的卖家在规则变动时几乎没有缓冲。
\end{insight}

\section{Temu 模式拆解}
\label{sec:plat-temu}

\subsection{从全托管到半托管与本地仓的演进}

Temu 于 2022 年 9 月在美国市场上线，公开信息普遍将其归为拼多多体系下的跨境平台
（平台背景与股权结构请以官方披露为准，见 \cite{temu_about}）。它最早的核心形态是典型的\kw{全托管}：
卖家只需要把货送到平台在国内的集货仓，
选品、定价、投流、跨境干线、海外履约、售后全部由平台完成。
这套模式的价值在于\kw{把跨境经营的门槛降到了「会做货」这一个维度}——
一个只会接单排产的工厂，不需要懂英语、不需要懂广告、
不需要在海外建仓，就能把货卖到几十个国家。

但全托管有一个内生的矛盾：\kw{平台的履约成本是刚性的，而低价是战略性的}。
空运小包在早期能够支撑「又快又便宜」的心智，
但随着单量上升与品类从轻小件扩展到中大件，空运的单位成本变得难以承受。
于是演进方向必然是\kw{把物流责任向卖家转移，同时把海外库存的效率优势拿回来}。
半托管（公开报道中多在 2024 年前后于美国等市场上线，具体时间与适用范围以官方公告为准）
就是这一转移的产物：卖家自己把货备到海外仓并从本地发货，
平台继续掌握定价与流量。
之后又出现「本地仓 / 本地对本地」形态，
即由本地卖家（或在当地有主体与库存的卖家）供货与履约，
进一步缩短时效并规避部分跨境合规摩擦。

\begin{figure}[htbp]
\centering
\begin{tikzpicture}[
  node distance=5mm,
  st/.style={draw=NavyMid, fill=BoxBG, rounded corners=2pt, minimum height=8mm,
             inner sep=3pt, text width=2.5cm, align=center, font=\sffamily\tiny},
  dec/.style={draw=AmberDeep, fill=white, diamond, aspect=2.2, inner sep=1pt,
              text width=1.5cm, align=center, font=\sffamily\tiny},
  bad/.style={draw=CrimsonDeep, fill=white, rounded corners=2pt, inner sep=3pt,
              text width=3.4cm, align=center, font=\sffamily\tiny},
  note/.style={draw=RuleGray, fill=BoxBGAlt, rounded corners=2pt, inner sep=4pt,
               text width=12.4cm, align=left, font=\sffamily\tiny}
]
\node[st] (a) at (0,0) {卖家提报\\图片与样品};
\node[st] (b) at (3.4,0) {买手选品\\同款比价};
\node[dec] (c) at (6.8,0) {核价\\通过?};
\node[st] (d) at (10.2,0) {上架测款\\小批量试单};
\node[st] (e) at (10.2,-2.1) {国内仓集货};
\node[st] (f) at (6.8,-2.1) {跨境干线\\与清关};
\node[st] (g) at (3.4,-2.1) {海外履约\\与尾程};
\node[st] (h) at (0,-2.1) {账期结算\\与售后追偿};
\node[bad] (x) at (6.8,1.9) {未通过：改款、降价\\或直接放弃};
\node[note] at (5.1,-3.8)
  {核价是全流程的\emph{唯一闸门}：报价一旦被判定偏高，
   后续所有环节都不会启动。因此卖家的优化重点不是「上架后如何提价」，
   而是「在报价环节如何用数据证明自己的价格是合理的」。};
\draw[-{Stealth}, thick, SlateGray] (a) -- (b);
\draw[-{Stealth}, thick, SlateGray] (b) -- (c);
\draw[-{Stealth}, thick, TealDeep] (c) -- (d) node[midway, above=1pt, font=\sffamily\tiny] {是};
\draw[-{Stealth}, thick, CrimsonDeep] (c) -- (x) node[midway, left=1pt, font=\sffamily\tiny] {否};
\draw[-{Stealth}, thick, SlateGray] (d) -- (e);
\draw[-{Stealth}, thick, SlateGray] (e) -- (f);
\draw[-{Stealth}, thick, SlateGray] (f) -- (g);
\draw[-{Stealth}, thick, SlateGray] (g) -- (h);
\end{tikzpicture}
\caption{Temu 模式的核价与备货流程：核价作为唯一准入闸门}
\label{fig:temu-flow}
\end{figure}

\subsection{核价机制：比价与压价的逻辑}

\term{核价}{Price Verification} 是全托管体系里最重要的一个概念，
也是卖家抱怨最多、误解最深的一个环节。
把公开可见的现象抽象出来，核价大致同时满足三条约束：
\kw{同款比价}（同款或高度相似款在平台内的历史最低成交价）、
\kw{站外比价}（同款在主流电商平台的可比价格）、
以及\kw{目标毛利倒推}（平台要求该品类贡献的毛利水平）。
把它写成一个可行域，卖家的报价 $p_{w}$ 需要满足
\begin{equation}
p_{w} \;\le\; p_{w}^{\max}
= \frac{p_{r}^{\text{目标}} - c_{\text{履约}} - c_{\text{营销}} - \pi_{\text{平台目标}}}{\gamma},
\qquad
p_{r}^{\text{目标}} = \min\Big( p_{r}^{\text{锚}},\;\; p_{r}^{\text{同类最低}} - \delta \Big),
\label{eq:temu-quote}
\end{equation}
其中 $p_{r}^{\text{目标}}$ 是平台为该商品设定的目标零售价，
$p_{r}^{\text{锚}}$ 是平台根据站外同款价格确定的锚点，
$p_{r}^{\text{同类最低}}$ 是平台内同类商品的当前最低零售价，
$\delta \ge 0$ 是平台要求的「显著更便宜」的价差（这是低价心智的工程化实现），
$\gamma$ 是从供货价到零售价的\kw{综合加价倍率}（隐含覆盖了平台毛利），
$c_{\text{履约}}$ 与 $c_{\text{营销}}$ 分别是平台侧的履约与获客成本。
这个式子的三个推论值得记住：

\begin{itemize}
  \item \textbf{核价是「倒推」而不是「议价」。} 卖家的成本结构与平台的定价逻辑几乎没有关系；
        平台先定零售价，再倒推它能接受的供货价。
        所以「我的成本就是这么高」在核价环节几乎没有说服力，
        能起作用的是\kw{工艺改进、包装简化、材料替代、产线分摊}这类把成本真的压下来的动作。
  \item \textbf{$\delta$ 是低价的发动机，也是生态的脆弱点。}
        只要有新卖家愿意用更低价格试单，$p_{r}^{\text{同类最低}}$ 就会被刷新，
        进而压低整条品类线的 $p_{w}^{\max}$。
        这是\kw{单向棘轮}：价格能下不能上。
        理解这一点，就能理解为什么全托管卖家的长期毛利预期必须按「逐步下行」来建模。
  \item \textbf{报价的真实下限由罚款与退货决定，而不是由成本决定。}
        见 \cref{eq:temu-cost}：因为赔付与退货损耗按金额比例发生，
        边际上会随着报价升高而放大。
\end{itemize}

把成本核算写成「可接报价」的形式，卖家在报价前应该自己先算清楚一条红线：
\begin{equation}
p_{w}^{\min}
= \frac{c_{\text{物料}} + c_{\text{加工}} + c_{\text{包装}}
+ c_{\text{国内干线}} + c_{\text{合规检测}} + F_{\text{固定}}/Q}
{1 - \lambda_{\text{赔付}} - \kappa_{\text{退货}}\, \theta_{\text{退货}} - m_{\text{目标}}} ,
\label{eq:temu-cost}
\end{equation}
其中 $\theta_{\text{退货}}$ 是退货造成的货值损失比例（能二次销售的退货只损失运费与折价，
不能二次销售的损失接近 100\%），$m_{\text{目标}}$ 是卖家自己要求的毛利率。
\cref{eq:temu-cost} 与常见的「成本加成」写法差在分母上：
\textbf{所有按金额比例发生的损失都必须放在分母里，而不是当作固定费用减在分子上}。
忽略这一点，是新手报价最常见的系统性错误——
一个 $\lambda_{\text{赔付}} = 3\%$、$\kappa_{\text{退货}} = 8\%$、
$\theta_{\text{退货}} = 0.5$ 的品类，等效损失约 7\%，
在 12\% 的目标毛利下会把报价低估将近 8\%。
这里的比例都是\kw{示范量级}，实际值因品类、季节与平台规则而异。

\subsection{买手制与选品流程}

全托管的选品不是「卖家上架、用户选择」，而是\kw{买手（品类采购）主导的招商制}。
流程通常是：卖家提交商品信息与样品，
买手在同款池中做横向比较（价格、评价、供应稳定性、交期），
通过核价后进入小批量试单，试单数据达标才会进入正式备货与推广。
这套机制有三个后果，卖家必须提前接受：

\begin{itemize}
  \item \textbf{选品是「被选中」而不是「被算法发现」。} 平台侧的推广资源是稀缺的，
        买手的注意力是更稀缺的资源。
        因此供应商的「配合度」（能否快速改款、能否配合特殊包装、能否接受小单）
        在长期合作中的权重，往往高于单纯的报价低一两毛钱。
  \item \textbf{同款竞争是常态，且难以用品牌规避。} 在核价体系里，
        同款就是同款，品牌故事的权重有限。
        卖家的护城河只能来自\kw{别人做不出或做不快}的工艺，
        比如特殊模具、独家面料、专利结构、难以复制的品控稳定性。
  \item \textbf{样品与打样成本由卖家前置承担。} 一个季度提报几十个款、
        只有少数通过核价、其中又只有一部分能跑量，是很正常的漏斗。
        把打样成本按「预期通过率」摊进成功款的成本里，
        是报价时必须做的会计处理（也就是 \cref{eq:temu-cost} 中的 $F_{\text{固定}}/Q$）。
\end{itemize}

\subsection{JIT 与备货模式}

\term{JIT}{Just-In-Time} 在跨境语境下指的是一种\kw{「先有订单、后发货」}的履约安排：
卖家不提前把货压到平台仓，而是在消费者下单后才把商品发往平台的国内集货仓，
由平台完成后续跨境履约。与之相对的是\kw{备货模式}：卖家提前把预测销量对应的库存送仓，
由平台从仓内直接发货。两者的取舍可以用一张表说清。

\begin{table}[htbp]
\centering
\caption{JIT 与备货模式的取舍}
\label{tab:temu-replenish}
\small
\begin{tabularx}{\textwidth}{L{2.5cm} Y Y}
\toprule
\tabhead{对比项} & \tabhead{JIT（出单后发货）} & \tabhead{备货（提前送仓）} \\
\midrule
资金占用 & 低，几乎按单生产 & 高，需为预测销量垫付 \\
\addlinespace
滞销风险 & 低，卖不动就停产 & 高，滞销库存需退回或处理 \\
\addlinespace
时效与流量 & 履约时效较长，可能影响前台权重与转化 & 时效短，通常更容易获得流量倾斜 \\
\addlinespace
缺货风险 & 高，产能波动直接变成缺货与罚款 & 低（若预测准确） \\
\addlinespace
适用场景 & 新款测款期、长尾款、定制款 & 爆款、稳定动销款、旺季前铺货 \\
\addlinespace
核心能力 & 快速响应的小批量产线 & 需求预测与库存计划 \\
\bottomrule
\end{tabularx}
\end{table}

这里有一个容易忽略的细节：\textbf{JIT 并没有消灭库存风险，
它只是把风险从「成品库存」转移到了「产能预留」}。
为了做到「出单后 24 到 48 小时能发货」，
工厂必须保持一定的空闲产能、常备原材料、
或者让上游面料商为自己留货——这些都是有成本的。
如果这条产线的空闲成本高于滞销库存的期望损失，
那么 JIT 反而更贵。判断标准很简单：
把「产能预留成本」与「滞销期望损失」写成同一个量纲比较即可，
滞销期望损失可以用 \cref{eq:shein-batch} 的同款思路计算。
这里还有一个容易被忽视的放大器：
备货模式下的需求预测本身会被\kw{牛鞭效应}扭曲
（需求信息在供应链上逐级放大与滞后，见 \cite{lee1997bullwhip}），
越是多级分销与多批补货的结构，预测误差越大；
因此备货模式对预测能力的要求远高于直觉，
标准的安全库存与再订货点方法见 \cite{silver1998inventory,hyndman2021forecasting}，
供应链层面的整体设计见 \cite{chopra2019supply}，
这些内容会在 \cref{ch:supply} 中系统展开。

\subsection{罚款与售后规则：全托管的隐性成本}

平台对供货商的约束手段，核心是\kw{赔付与罚款}。
公开可见的常见事由包括：品控不合格、到货超时、缺货（未按承诺数量交付）、
标签与包装不合规、虚假宣传、侵权、以及售后质量问题（退货、退款、客诉）。
具体条款、计算方式与豁免条件\kw{以各平台官方最新规则为准}，
但它们的\kw{数学结构}是稳定的：罚款通常由「按单固定罚款」、
「按批次固定罚款」与「按金额比例罚款」三种形式组合而成。
写成对单件利润的侵蚀：
\begin{equation}
\Delta \pi
= -\Big( f_{\text{单}} + \frac{F_{\text{批}}}{Q}
+ \lambda_{\text{比例}}\, p_{w} \Big),
\qquad
\frac{\partial \pi}{\partial \lambda_{\text{比例}}} = -\,p_{w} .
\label{eq:temu-penalty}
\end{equation}
\cref{eq:temu-penalty} 传递了三个可以立刻用于决策的结论：

\begin{enumerate}
  \item \textbf{罚款率与报价同向放大。} 罚的是金额比例，
        所以报价越高，同样的违规率造成的绝对损失越大；
        这进一步压缩了「靠涨价改善利润」的空间（而涨价本身还要过核价这一关）。
  \item \textbf{批次数越多，按批罚款的总量越大。}
        一个把 10000 件分成 20 个批次交付的卖家，
        在按批罚款的规则下暴露的风险是单批次交付的 20 倍。
        这就是\kw{合并批次、提高单批完整率}的财务动机。
  \item \textbf{固定罚款对小单价的伤害是致命的。}
        如果 $f_{\text{单}} = 20$ 元而单件利润只有 5 元，
        那么一次罚款需要 4 件的利润来弥补。
        低价小件品类尤其要警惕这一点：\kw{单价越低，容错空间越小}。
\end{enumerate}

\begin{pitfall}[「全托管看起来没风险」——备货滞销与罚款的隐性成本]
\textbf{现象}：一个做家居小件的工厂听说全托管「不用管运营、不用管物流、不用管售后」，
一次性备了三个柜的货送到平台国内仓，前两个月动销不错，
第三个月被通知部分 SKU 需要降价清仓，同时因两批货标签不合规被罚款。
半年下来财务核算，账面供货收入看起来不错，但净利润接近零甚至为负。\\
\textbf{根因}：这里有三个被系统性低估的成本项。
(1) \textbf{滞销不是「退货」，而是「退回加折价」}：卖不动的库存退回时的运费、
以及二次销售的价值折损，通常由卖家承担；
(2) \textbf{罚款是按金额比例发生的}（\cref{eq:temu-penalty}），
一旦品控或标签出问题，损失会随供货额同步放大；
(3) \textbf{账期把利润变成了应收款}（\cref{eq:fm-cashgap}），
财务上「赚到钱」和「拿到钱」之间隔着一个现金缺口，
如果这个缺口靠借贷填补，利息又要再吃掉一层利润。\\
\textbf{排查方法}：
(1) 把「单件净利」定义为\textbf{扣除罚款、退货货值损失与资金成本之后}的口径，
并在系统里把它作为唯一的考核口径，禁止用「供货价减成本」汇报业绩；
(2) 对每个 SKU 设定\kw{滞销止损线}（例如上架 30 天动销率低于某一阈值即触发降价或停产），
把它写进运营流程而不是靠感觉；
(3) 罚款与退货数据按月做归因，区分「规则理解问题」「工艺问题」「包装问题」三类，
分别对应培训、工艺改进与包材修改，而不是笼统地归为「平台太苛刻」；
(4) 在报价阶段就把 \cref{eq:temu-cost} 中的 $\lambda_{\text{赔付}}$ 与
$\kappa_{\text{退货}}$ 用历史数据标定，而不是填零。
\end{pitfall}

\subsection{增长打法：低价、社交裂变与超级碗投放}

从公开报道与平台自身的营销动作可以观察到一个相对清晰的增长组合：
\kw{极致低价建立比价心智} 加上 \kw{社交裂变降低获客边际成本}，
再用\kw{超级碗等大曝光事件建立「便宜」的全民认知}。
这套打法的经济学逻辑可以用一个简单的分解来理解：
把获客成本摊到订单上，\metric{CAC} 由自然流量、裂变流量与付费流量三部分混合而成；
当裂变与自然流量占比提高时，混合 \metric{CAC} 会显著低于纯付费投放
（裂变与推荐机制的增长逻辑可参考 \cite{ellis2021growth,ries2011lean}）。
大曝光投放的作用不只是当期转化，
而是把「这个平台便宜」变成消费者的\kw{默认假设}，
从而让后续每一次搜索与点击的转化率都更高——
这是一种把营销费用转化为\kw{长期认知资产}的支出。
公开报道中关于单次超级碗广告位价格等数字（有报道给出千万美元量级的口径）
存在不确定性，且不同年份、不同时段差异很大，
因此这里只作方向性描述，不作为任何测算依据，
相关费用的官方口径请以平台与广告主披露为准。

\begin{casebox}[案例：一个工厂型卖家的双平台利润结构对比（示范量级）]
\textbf{说明}：以下全部为便于说明结构的\textbf{示范计算}，
数字为演示量级，不代表任何真实企业、真实平台或真实品类的经营数据。
所有费率、佣金与账期请以平台官方最新规则为准。

设某工厂生产一款家居小件，出厂完全成本（物料加加工加包装）为 12 元。
该工厂同时在两个渠道销售：全托管平台与自主运营的内容电商。

% tcolorbox 内禁止浮动体，这里用 center 加 captionof 生成不浮动表格。
\begin{center}
\small
\begin{tabularx}{\linewidth}{L{2.6cm} L{3.1cm} L{3.4cm} Y}
\toprule
\tabhead{项目} & \tabhead{全托管渠道} & \tabhead{内容电商自运营} & \tabhead{说明} \\
\midrule
出厂成本 & 12.0 元 & 12.0 元 & 同一条产线 \\
\addlinespace
收入口径 & 供货价 18.0 元 & 零售价约 20 美元 & 币种与价格带差异极大 \\
\addlinespace
平台侧扣除 & 罚款与赔付约 3\% & 佣金与支付费约 8\% & 以官方规则为准 \\
\addlinespace
物流与关税 & 国内段约 0.6 元 & 头程、关税、仓配合计约 25\% & 自运营显著更高 \\
\addlinespace
流量成本 & 近乎为零（平台承担） & 广告约 15\% 至 25\% & 自运营最大变量 \\
\addlinespace
退货与售后 & 退货率约 6\% & 退货率约 8\% & 服饰类通常更高 \\
\addlinespace
单件净利 & 约 3.5 至 4.5 元 & 约 2.5 至 4.0 美元 & 折人民币后自运营更高 \\
\addlinespace
资金占用 & 低（送到国内仓即可） & 高（海外备货加广告预付） & 决定谁能规模化 \\
\bottomrule
\end{tabularx}
\captionof{table}{工厂型卖家双平台利润结构对比（示范量级，非真实数据）}
\label{tab:case-dual}
\end{center}

\textbf{结论与取舍}：全托管渠道的优势是\kw{资金占用低、起量快、不需要团队}，
劣势是\kw{单件利润薄且价格只能下行}；
自运营渠道的优势是\kw{单件利润高、可累积品牌与用户}，
劣势是\kw{需要备货到海外、需要投放与内容团队、波动大}。
理性的做法通常不是二选一，而是用全托管渠道\kw{消化产能与现金流}，
用自运营渠道\kw{试品牌与高毛利款}，
并把全托管赚到的钱按固定比例投入到自运营的库存与内容建设上——
本质上是把一个低风险低毛利的业务，当作高风险高毛利业务的融资来源。
\end{casebox}

\subsection{卖家的真实利润率与风险}

把前面的内容收拢，全托管卖家面对的风险可以归为四类，
它们恰好对应 \cref{eq:fm-profit} 中的四个负项：

\begin{itemize}
  \item \textbf{被压价（单价风险）}：核价机制的单向棘轮结构，
        使得长期毛利预期必须按逐步下行建模。
        应对方式是\kw{持续做成本工程}，而不是指望一次谈判定终身。
  \item \textbf{库存滞销（批量风险）}：动销不及预期时，
        滞销库存要么退回、要么降价，两者都损失利润。
        应对方式是控制单批规模、用 JIT 承接不确定性、设定止损线。
  \item \textbf{罚款与赔付（规则风险）}：随金额比例放大。
        应对方式是把品控、标签、包装、时效做成\kw{可审计的流程}，
        而不是依赖一线工人的个人经验。
  \item \textbf{资金占用（账期风险）}：利润变成应收款，缺口靠融资填补。
        应对方式是争取更短的账期或更长的上游账期，
        并把资金成本显式计入报价（这是很多卖家在报价时忘掉的一项）。
\end{itemize}

\begin{formulacard}[公式速记：全托管卖家的单件净利]
\begin{equation}
\pi_{\text{净}}
= \underbrace{p_{w} - c_{\text{全成本}}}_{\text{毛差}}
- \underbrace{\lambda_{\text{赔付}}\, p_{w}}_{\text{规则风险}}
- \underbrace{\kappa_{\text{退货}}\, \theta_{\text{退货}}\, c_{\text{全成本}}}_{\text{退货风险}}
- \underbrace{\frac{\text{年资金成本} \times \text{现金缺口}}{\text{年销量}}}_{\text{账期成本}} .
\label{eq:fm-net}
\end{equation}
记住三句话：\textbf{毛差决定天花板，赔付率决定稳定性，
账期决定你明年还在不在牌桌上}。
\end{formulacard}

\section{SHEIN 模式拆解：小单快反与柔性供应链}
\label{sec:plat-shein}

\subsection{小单快反：首单、测款与返单}

SHEIN 最被反复讨论的能力是\term{小单快反}{Small-Batch Rapid Response}：
每个新款先下\kw{极小批量的首单}（公开报道中常提到百件量级的口径，
具体数量因品类而异，见 \cite{shein_about}），
上架后根据\kw{真实销售数据}在极短时间内决定返单还是停产。
它看起来只是一个「下小单」的操作习惯，实际上是一整套
\kw{把库存风险转化为信息成本}的机制设计。

把这个决策写成期望利润的形式，一个批次的期望利润为
\begin{equation}
\E[\pi_{\text{批次}}]
= Q\Big[\, s\,(p_{r} - c_{\text{变动}}) - (1-s)\,\theta_{\text{滞销}}\, c_{\text{变动}} \,\Big]
- c_{\text{测款}} - F_{\text{打样}} ,
\label{eq:shein-batch}
\end{equation}
其中 $Q$ 是批次数量，$s \in [0,1]$ 是该批次的\kw{售罄率}，
$p_{r}$ 是零售价，$c_{\text{变动}}$ 是单件变动成本，
$\theta_{\text{滞销}} \in [0,1]$ 是滞销库存的价值损失比例，
$c_{\text{测款}}$ 与 $F_{\text{打样}}$ 是每个 SKU 的前置固定投入。
令 \cref{eq:shein-batch} 为零，可以解出\kw{返单的临界售罄率}：
\begin{equation}
s^{*}
= \frac{\theta_{\text{滞销}}\, c_{\text{变动}} + \dfrac{c_{\text{测款}} + F_{\text{打样}}}{Q}}
{p_{r} - c_{\text{变动}} + \theta_{\text{滞销}}\, c_{\text{变动}}} .
\label{eq:shein-threshold}
\end{equation}
\cref{eq:shein-threshold} 揭示了小单快反最关键的两个设计参数：

\begin{itemize}
  \item \textbf{$Q$ 越小，$s^{*}$ 越高。} 因为前置固定成本被摊到更少件数上，
        单件必须更成功才回本。
        这就是小单的代价：\kw{用更高的单件成本换取更低的滞销风险}。
        首单到底下 100 件还是 300 件，本质上是在这两者之间找平衡点。
  \item \textbf{$p_{r} - c_{\text{变动}}$ 越大，$s^{*}$ 越低。}
        高毛利款式能够承受更高的测款失败率，
        这解释了为什么快时尚品类特别适合小单快反——
        它的零售加价倍率足够高，容得下大量失败款。
        \textbf{低毛利标品几乎无法照搬这套打法}，
        因为单件毛差太小，任何固定成本摊销都会让临界售罄率高到不可实现。
\end{itemize}

\begin{insight}[小单快反真正卖的是「信息」，不是「衣服」]
把 \cref{eq:shein-batch,eq:shein-threshold} 放在一起看，
会发现小单快反的收益并不主要来自「少生产」，
而是来自\kw{用真实销售数据替代需求预测}。
传统模式必须先预测再生产，预测误差直接变成库存损失；
小单快反先把不确定性隔离在一个很小的 $Q$ 上，
用 $Q$ 的销售结果作为\kw{无偏的需求信号}，再决定是否放大。
换句话说，它把供应链的决策顺序从「预测、生产、销售」
改成了「小规模生产、观察、放量、销售」。
这个改动之所以难复制，不是因为它需要复杂的算法，
而是因为它需要\kw{整条供应链愿意配合极小批量}——
面料商愿意拆卷、工厂愿意频繁换线、物流愿意高频小批集货，
这些是产业带密度与长期关系带来的能力，不是买一套系统就能获得的。
\end{insight}

\subsection{柔性供应链：产业带集群与打样到上架}

SHEIN 的供应链公开信息中，最常被提及的是\kw{广州番禺一带的服装产业带集群}：
在很小的地理半径内聚集了大量面料商、辅料商、打版房、裁床、车缝厂与后整厂，
使得「设计、打样、上架」的周期可以压缩到公开报道所说的\kw{最短 3 到 7 天}
（具体周期随品类与季节变化，属公开报道口径，存在不确定性）。
产业带密度的价值可以用一句话概括：
\kw{物理距离就是交付周期，而交付周期就是库存}。

\begin{table}[htbp]
\centering
\caption{传统大批量模式与小单快反模式的对比}
\label{tab:shein-compare}
\small
\begin{tabularx}{\textwidth}{L{2.6cm} Y Y}
\toprule
\tabhead{对比项} & \tabhead{传统大批量（预测驱动）} & \tabhead{小单快反（数据驱动）} \\
\midrule
决策依据 & 历史销量、买手经验与趋势判断 & 首单真实动销数据 \\
\addlinespace
单 SKU 首单量 & 千件以上 & 百件量级（品类相关） \\
\addlinespace
补货周期 & 数周至数月 & 数天至两周 \\
\addlinespace
单件成本 & 低（规模效应） & 高（换线、小批采购溢价） \\
\addlinespace
库存周转 & 数十天至上百天 & 公开报道口径为个位数天 \\
\addlinespace
失败成本 & 高，滞销库存难以消化 & 低，损失被限制在首单 \\
\addlinespace
组织要求 & 计划与采购能力强 & 买手、打版、产能调度极强 \\
\addlinespace
适配品类 & 标准品、季节性弱、需求稳定 & 时尚属性强、需求波动大、加价倍率高 \\
\bottomrule
\end{tabularx}
\end{table}

\subsection{库存周转的经济学意义}

库存周转天数是这套模式里最容易被低估的财务变量。
把周转与资金效率写成公式：
\begin{equation}
\text{年周转次数} = \frac{365}{D_{\text{库存}} + D_{\text{在产}}},
\qquad
\ROI_{\text{资金}}
= \frac{\pi_{\text{单件}} \cdot Q_{\text{年销}}}{W_{\text{平均占用}}}
= \frac{\pi_{\text{单件}}}{c_{\text{单件}}} \times \text{年周转次数} ,
\label{eq:shein-turn}
\end{equation}
其中 $W_{\text{平均占用}}$ 是平均占用资金。
\cref{eq:shein-turn} 的第二个等号是关键：
\kw{资金回报率等于单件毛利率乘以年周转次数}。
这意味着一个「单件毛利率 10\%、年周转 12 次」的生意，
资金回报率约 120\%，远高于「单件毛利率 40\%、年周转 1.5 次」的生意（约 60\%）。
\textbf{库存周转不是运营指标，它是财务指标}——
它和毛利率是同一个公式里的两个因子，任一个改善都能线形提升资金回报。

\begin{keypoint}[周转率是低毛利业务的唯一出路]
当单件毛利被平台核价压到个位数时，卖家只有两个方向可选：
\kw{提高周转}（更快卖、更少压货、更短账期），
或者\kw{提高单件毛利}（做差异化、做品牌、换品类）。
如果两条都不做，那么名义上「还有 5\% 毛利」的生意，
在计入资金成本、罚款与滞销损失之后一定是亏的。
把 \cref{eq:shein-turn} 打印出来贴在办公室，
比任何「精细化运营」的口号都更能指导决策：
\textbf{要么让钱转得更快，要么让每转一圈赚得更多。}
\end{keypoint}

\subsection{数字化选品与测款机制}

柔性供应链解决的是「能不能快速生产」，数字化选品解决的是「生产什么」。
公开信息中可见的做法包括：抓取站内搜索与浏览行为、
快速识别上升趋势的关键词与款式元素（领型、袖型、图案、色彩），
把趋势信号转化为设计需求，再用\kw{小批量测款}在真实流量中验证。
这里有一个很重要的工程细节：
\kw{测款不是均匀铺量}，而是把有限的曝光资源按预期潜力分配，
让更有希望的款更快获得统计上可判断的数据量。
这与推荐系统里的\kw{探索与利用}问题是同一个数学结构，
相关方法可回看 \cref{ch:recall}。

测款要产生可靠的决策，前提是\kw{数据量达到可判断的水平}。
如果用 $n$ 次曝光的转化率 $s$ 作为判断依据，
其标准误约为 $\sqrt{s(1-s)/n}$，
在 $s = 0.02$、$n = 500$ 时标准误约 0.0063，
也就是 0.63 个百分点，与转化率本身的量级相当。
这意味着\kw{几百次曝光的测款结果几乎没有区分度}，
用它做「好款或坏款」的二元判断很容易被噪声主导。
正确的做法是：要么把测款曝光量提到千次以上，
要么改用排序（比较多个款的相对表现）而不是绝对阈值来决策。
这个细节是很多团队「测款测不准」的真正原因。

\subsection{平台化转型：SHEIN Marketplace 的规则}

从 2023 年前后开始，SHEIN 逐步从纯自营走向\kw{平台化}，
面向第三方卖家开放 Marketplace 招商（公开报道中首站为北美等市场，
后续逐步扩展；具体开放站点、类目与费率\kw{以官方最新规则为准}）。
这一步的商业含义是：把已经建成的\kw{流量与履约基础设施}
开放给外部供给，从而在保持自营供应链优势的同时，
快速扩充品类宽度（尤其是服装之外的品类）。

\begin{table}[htbp]
\centering
\caption{SHEIN Marketplace 模式卖家通常需要关注的规则维度（以官方最新规则为准）}
\label{tab:shein-marketplace}
\small
\begin{tabularx}{\textwidth}{L{2.6cm} L{4.8cm} Y}
\toprule
\tabhead{维度} & \tabhead{公开规则口径下的关注点} & \tabhead{对卖家的实际影响} \\
\midrule
定价与核价 & 平台对价格竞争力有要求，可能设置价格对齐机制 & 定价空间受限，需提前测算让价后的利润 \\
\addlinespace
佣金与费用 & 按类目收取佣金，另有支付、履约等费用 & 直接决定毛利结构，需逐类目核算 \\
\addlinespace
履约方式 & 可能要求卖家发货至平台仓，或允许本地发货 & 决定资金占用与时效表现 \\
\addlinespace
退货与售后 & 平台规则下卖家承担相应售后成本 & 高退货类目需提高退货率假设 \\
\addlinespace
合规与认证 & 目的国产品合规、标签与税务登记要求 & 准入门槛，未准备好会直接卡在入驻环节 \\
\bottomrule
\end{tabularx}
\end{table}

\begin{pitfall}[把「平台招商」当成「免费流量」]
\textbf{现象}：看到平台招商就立刻入驻、大量铺货，
半年后发现「有流量但没有利润」，甚至因为价格对齐机制被迫降价到亏损。\\
\textbf{根因}：平台招商的动机是\kw{扩充品类宽度}，不是\kw{分配免费流量}。
在成熟平台的流量分配里，新卖家能拿到的初始资源通常只够完成「测款」，
真正的持续性流量仍要靠价格竞争力、履约表现与商品质量分去换。
把招商政策里的「流量扶持」理解成长期承诺，是最常见的预期错配。\\
\textbf{排查方法}：
(1) 入驻前先做\kw{单位经济学测算}——把佣金、履约、退货与可能的降价全部计入，
算出「在最坏价格情形下单件是否为正」；
(2) 用小规模测款验证真实流量获取成本，而不是用招商材料里的案例数据；
(3) 明确「如果平台要求价格对齐到某个水平，我是否还愿意做」，
把这条作为硬约束写进选品标准，而不是事后被动接受。
\end{pitfall}

\section{TikTok Shop 模式拆解：内容与交易的一体化}
\label{sec:plat-tiktok}

\subsection{内容--交易一体：从短视频到小店成交}

TikTok Shop 与前面两个平台的根本区别在于：
\kw{它不是先有需求再成交，而是先创造需求再成交}。
Temu 的用户带着「找便宜」的意图来，SHEIN 的用户带着「看新款」的意图来，
而短视频平台的用户是来\kw{消费内容}的——
成交必须建立在内容本身足够有吸引力的前提上（平台定位见 \cite{tiktokshop_about}）。
这条路径可以拆成一个漏斗：
\kw{内容曝光} 到 \kw{视频或直播间的进入与停留}，
再到 \kw{商品卡点击}、\kw{小店下单}、\kw{履约} 与 \kw{复购}。
把它写成 GMV 的分解式：
\begin{equation}
\GMV
= N_{\text{曝光}} \cdot r_{\text{进入}} \cdot \CTR_{\text{商品点击}}
\cdot \CVR_{\text{下单}} \cdot \AOV \cdot (1 - \rho_{\text{退款}}) .
\label{eq:tt-live-gmv}
\end{equation}
\cref{eq:tt-live-gmv} 的价值在于把「人货场」这个行业说法翻译成了可优化变量：
\kw{场}对应 $r_{\text{进入}}$（封面、标题、直播贴片、场景布置），
\kw{人}对应 $\CTR_{\text{商品点击}}$ 与 $\CVR_{\text{下单}}$（主播讲解、信任、演示），
\kw{货}对应 $\AOV$ 与 $\rho_{\text{退款}}$（商品力、价格、质量与描述一致性）。
任何一场直播的复盘，都应该先把这六个因子跑一遍，
而不是笼统地说「今天流量不好」。

\begin{figure}[htbp]
\centering
\begin{tikzpicture}[
  every node/.style={align=left},
  st/.style={draw=NavyMid, fill=BoxBG, rounded corners=2pt, minimum height=7mm,
             inner sep=4pt, text width=4.6cm, font=\sffamily\tiny},
  mt/.style={draw=TealDeep, fill=BoxBGAlt, rounded corners=2pt, minimum height=7mm,
             inner sep=4pt, text width=6.4cm, font=\sffamily\tiny},
  note/.style={draw=RuleGray, fill=white, rounded corners=2pt, inner sep=4pt,
               text width=12.4cm, align=left, font=\sffamily\tiny}
]
\node[st] (a) at (0,0)      {1. 内容曝光：短视频与直播推荐};
\node[st] (b) at (0,-1.15)  {2. 进入与停留：封面、标题、场景};
\node[st] (c) at (0,-2.30)  {3. 商品卡点击：讲解、演示、信任};
\node[st] (d) at (0,-3.45)  {4. 小店下单：价格、库存、优惠};
\node[st] (e) at (0,-4.60)  {5. 履约与售后：FBT 或自发货};
\node[st] (f) at (0,-5.75)  {6. 复购与达人关系沉淀};
\node[mt] (ma) at (5.5,0)      {$N_{\text{曝光}}$：内容产能、发布频率、账号权重};
\node[mt] (mb) at (5.5,-1.15)  {$r_{\text{进入}}$：封面点击率、贴片、开播时段};
\node[mt] (mc) at (5.5,-2.30)  {$\CTR_{\text{商品点击}}$：话术、场景、痛点演示};
\node[mt] (md) at (5.5,-3.45)  {$\CVR_{\text{下单}}$ 与 $\AOV$：价格力、组合、免运};
\node[mt] (me) at (5.5,-4.60)  {$\rho_{\text{退款}}$：描述一致性、质量、时效};
\node[mt] (mf) at (5.5,-5.75)  {$\LTV$：复购率、达人复投、粉丝沉淀};
\draw[-{Stealth}, thick, SlateGray] (a) -- (b);
\draw[-{Stealth}, thick, SlateGray] (b) -- (c);
\draw[-{Stealth}, thick, SlateGray] (c) -- (d);
\draw[-{Stealth}, thick, SlateGray] (d) -- (e);
\draw[-{Stealth}, thick, SlateGray] (e) -- (f);
\end{tikzpicture}
\caption{TikTok Shop 的内容--交易漏斗与对应的可优化变量}
\label{fig:tt-funnel}
\end{figure}

\subsection{美区与东南亚：两套完全不同的市场结构}

同一个产品、同一套达人机制，在美区与东南亚的表现差异极大。
造成差异的根本原因不在平台规则，而在\kw{市场结构}：
客单价、物流基础设施、支付习惯、内容消费习惯与竞争密度都不同。

\begin{table}[htbp]
\centering
\caption{美区与东南亚市场结构差异对运营策略的影响（定性总结，量为示范量级）}
\label{tab:tt-region}
\small
\begin{tabularx}{\textwidth}{L{2.4cm} Y Y}
\toprule
\tabhead{维度} & \tabhead{美区} & \tabhead{东南亚} \\
\midrule
客单价 & 较高，可承受一定的履约与广告成本 & 较低，价格敏感度极高 \\
\addlinespace
履约时效 & 消费者对时效预期高，本地仓几乎必需 & 时效容忍度相对更高，跨境直发仍有一定空间 \\
\addlinespace
退货率 & 相对较高，退货政策宽松 & 相对较低，但货到付款场景存在拒收风险 \\
\addlinespace
内容生态 & 达人报价高、专业度高，头部集中 & 达人数量多、报价低，中腰部达人性价比高 \\
\addlinespace
合规与税务 & 产品合规、税务登记与平台资质要求严格 & 各国差异大，需逐国处理 \\
\addlinespace
政策风险 & 存在平台所有权与数据安全相关的立法与监管不确定性 & 存在各国监管与本地化经营要求的变化 \\
\bottomrule
\end{tabularx}
\end{table}

关于美区，必须单独说明\kw{政策风险}：
公开报道显示，围绕短视频平台的所有权与数据安全问题，
美国存在持续的立法与监管讨论（例如 2024 年前后出现的相关法案动向），
这类事件的走向具有高度不确定性，
可能通过股权结构变化、本地化运营要求或经营限制等方式影响平台生态。
对卖家而言，正确的应对不是预测结果，而是\kw{不把全部库存与现金流押在单一市场}，
具体讨论见 \cref{sec:plat-strategy}。
涉及数据跨境与内容合规的处理，必须回到 \cref{ch:cbec} 的框架逐项核对。

\subsection{全托管与自运营：两条入驻路径}

TikTok Shop 同时提供两条路径，它们分别对应前面定义的全托管与自主运营：

\begin{itemize}
  \item \textbf{全托管（平台运营、卖家供货）}：卖家按平台要求供货，
        定价、内容、达人合作、履约与售后主要由平台组织。
        卖家面对的是一套\kw{与全托管电商类似的核价与交付规则}，
        但平台侧的需求创造能力来自内容场，因此选品逻辑偏向
        \kw{「内容好讲、视觉冲击强、单价适中」的商品}。
  \item \textbf{自运营（商家自播加达人分销）}：卖家自己开店，
        自己发内容、自己开播、自己找达人、自己履约。
        定价权完整，但流量获取能力成为决定性门槛。
        这条路与买量逻辑高度重合，
        短视频平台的广告投放细节可参考 \cref{ch:meta}。
\end{itemize}

\subsection{Affiliate 达人分销机制：佣金、样品与计划类型}

\term{Affiliate}{达人分销} 是内容电商最核心的增长杠杆：
商家把商品加入联盟并设置佣金，达人（Creator）在内容里挂商品链接，
成交后按佣金率结算。它把\kw{广告费的前置现金支出}
变成了\kw{成交后的分佣支出}，显著降低了卖家的现金流压力，
但代价是把一部分毛利让给了渠道。

把单件利润写出来：
\begin{equation}
\pi_{\text{达人单}}
= p_{r}\,\big(1 - r_{\text{平台}} - r_{\text{达人}} - \rho_{\text{退款}}\big)
- c_{\text{货}} - c_{\text{履约}} - \frac{c_{\text{样品}}}{Q_{\text{达人}}} ,
\label{eq:tt-affiliate}
\end{equation}
其中 $r_{\text{平台}}$ 是平台佣金与支付费之和，
$r_{\text{达人}}$ 是达人佣金率，$\rho_{\text{退款}}$ 是退款率，
$c_{\text{履约}}$ 是尾程与仓储的摊薄成本，
$c_{\text{样品}}/Q_{\text{达人}}$ 是寄样成本（含未被使用、
未产生成交的样品）摊到该达人带来的成交件数上。
令 \cref{eq:tt-affiliate} 为零，得到\kw{可承受的最高佣金率}：
\begin{equation}
r_{\text{达人}}^{*}
= 1 - r_{\text{平台}} - \rho_{\text{退款}}
- \frac{c_{\text{货}} + c_{\text{履约}} + c_{\text{样品}}/Q_{\text{达人}}}{p_{r}} .
\label{eq:tt-breakeven}
\end{equation}
\cref{eq:tt-breakeven} 是达人合作谈判里最实用的一条公式。
举例来说（\kw{示范量级}），若 $r_{\text{平台}} = 0.08$、
$\rho_{\text{退款}} = 0.10$、货与履约成本占零售价的 35\%，
则 $r_{\text{达人}}^{*} \approx 0.47$，看上去空间很大；
但如果再加入 15\% 的内容投流摊销，
可承受佣金率立刻降到 0.32。
这解释了为什么「佣金给到 20\% 没人带、给到 30\% 自己亏」的矛盾会普遍出现——
问题往往不在佣金率本身，而在\kw{退款率与履约成本这两项被低估了}。

\begin{table}[htbp]
\centering
\caption{达人合作计划的两类常见形式对比（具体规则以官方最新文档为准）}
\label{tab:tt-plan}
\small
\begin{tabularx}{\textwidth}{L{2.5cm} Y Y}
\toprule
\tabhead{对比项} & \tabhead{Open Plan（公开计划）} & \tabhead{Targeted Plan（定向计划）} \\
\midrule
可见范围 & 全平台达人均可申请 & 仅对指定达人可见或可申请 \\
\addlinespace
佣金策略 & 通常为统一佣金率，覆盖面广 & 可对头部达人或特定类型达人给更高佣金 \\
\addlinespace
样品门槛 & 一般有申请与审核流程 & 通常可定向寄样，沟通成本更高 \\
\addlinespace
管理成本 & 低，靠规模自然筛选 & 高，需要一对一沟通与履约跟踪 \\
\addlinespace
适合阶段 & 冷启动、铺量测款 & 打爆款、维护核心达人关系 \\
\addlinespace
主要风险 & 佣金支出分散、转化不确定 & 佣金成本高，达人履约不稳定 \\
\bottomrule
\end{tabularx}
\end{table}

\begin{toolbox}[达人分销的样品投放清单]
寄样是达人分销里最容易亏钱的一环。把下面这套流程做成标准动作：
\begin{enumerate}
  \item \textbf{先筛选，再寄样}：按达人的\kw{历史带货类目相关性}、
        近期内容播放中位数、以及历史挂车成交表现做初筛，
        不要按粉丝数排序。
  \item \textbf{设定寄样预算上限}：按 \cref{eq:tt-breakeven} 反推——
        如果单件净利是 $\pi$、样品成本是 $c_{s}$，
        那么一个达人至少要带来 $\lceil c_{s}/\pi \rceil$ 件成交才回本。
        把这个数字作为寄样审批的门槛。
  \item \textbf{用阶梯佣金代替一口价}：
        对月成交低于阈值的达人给基础佣金，
        超过阈值后按阶梯上浮，把佣金支出与产出绑定。
  \item \textbf{跟踪三个指标而不是一个}：寄样后的\kw{内容发布率}、
        \kw{挂车率}与\kw{单达人产出}。
        只看成交额会漏掉「发了但不挂车」这个最常见的浪费。
  \item \textbf{沉淀可复用素材}：把表现好的达人内容拆成脚本结构，
        用于商家自播与自有素材，降低内容生产成本。
\end{enumerate}
\end{toolbox}

\subsection{联盟生态：机构、团长与达人分层}

随着达人数量增长，平台生态里自然演化出中间层：
\kw{联盟机构}（签约与管理达人）与\kw{团长}（组织达人分销的中间人）。
它们存在的经济理由和 \cref{ch:ad} 里讨论的中介层完全一样：
降低商家的\kw{达人筛选成本}与达人的\kw{选品成本}，
并承担一部分履约与结算风险。
商家在与机构合作时需要特别注意三点：

\begin{itemize}
  \item \textbf{分清「佣金」与「服务费」}：服务费通常是固定支出或按成交额抽成，
        与是否成交无关，会直接改变 \cref{eq:tt-breakeven} 的解。
  \item \textbf{明确数据归属}：达人带来的用户与内容素材的复用权，
        应在合同中写清；否则商家付出了佣金却沉淀不下任何资产。
  \item \textbf{警惕「纯拉群」型中间层}：不提供选品、素材或投放支持的中间层，
        对商家的边际价值极低，却要分走固定的佣金差额。
\end{itemize}

\subsection{直播间的「人货场」设计}

把「人货场」拆开，每一项都可以对应 \cref{eq:tt-live-gmv} 中的一个因子，
并且都有可执行的优化动作：

\begin{itemize}
  \item \textbf{场（对应 $r_{\text{进入}}$）}：直播间贴片、场景布置、
        背景文字、商品陈列顺序。核心目标是让\kw{划到的人愿意停 3 秒}。
        优化方法是做封面与贴片文案的对照实验，
        实验设计原则可参考 \cite{kohavi2020trustworthy}。
  \item \textbf{人（对应 $\CTR_{\text{商品点击}}$ 与 $\CVR_{\text{下单}}$）}：
        主播的讲解结构、演示方式、答疑与信任建设。
        有效的结构通常是\kw{痛点场景、解决方案演示、价格锚点、限时行动}，
        而不是从头到尾念参数。
  \item \textbf{货（对应 $\AOV$ 与 $\rho_{\text{退款}}$）}：
        价格带的分布、引流款与利润款的搭配、
        以及商品描述与实际的一致性。
        \kw{描述不一致是退款率飙升的第一大原因}，
        而退款率会同时伤害利润与账号权重，属于必须优先治理的项。
\end{itemize}

\subsection{FBT 与自发货}

\term{FBT}{Fulfilled by TikTok} 指由平台提供仓储与配送服务，
卖家提前把货备到平台的仓库，成交后由平台完成拣货、打包与配送；
\kw{自发货}则要求卖家（或第三方海外仓）自行完成本地履约。

\begin{table}[htbp]
\centering
\caption{FBT 与自发货的取舍（费用结构与规则以官方最新文档为准）}
\label{tab:tt-fulfil}
\small
\begin{tabularx}{\textwidth}{L{2.5cm} Y Y}
\toprule
\tabhead{对比项} & \tabhead{FBT（平台仓配）} & \tabhead{自发货或第三方海外仓} \\
\midrule
时效与转化 & 时效稳定，通常有利于转化与账号表现 & 取决于自选服务商，波动较大 \\
\addlinespace
费用结构 & 仓储费、操作费、配送费，另有长期仓储与移除费用 & 仓租与操作费可议价，配送费自谈 \\
\addlinespace
库存风险 & 滞销仍会产生仓储与移除成本 & 灵活度更高，可自行调拨与清仓 \\
\addlinespace
运营复杂度 & 低，与平台系统打通 & 高，需要对接仓储与订单系统 \\
\addlinespace
合规 & 由平台承担大部分消费者侧合规 & 卖家自行承担 \\
\addlinespace
适用阶段 & 爆款期、追求转化与稳定的阶段 & 多平台共用库存、需要灵活调拨的阶段 \\
\bottomrule
\end{tabularx}
\end{table}

\begin{insight}[FBT 的真正价值是「权重」，不是「便利」]
很多人把 FBT 当成「省事的外包」，这是低估了它。
在内容电商里，\kw{履约时效直接进入流量分配的考量}：
同样的内容转化能力，时效稳定的一方在同样的曝光下
能获得更高的成交概率与更低的退款率，
进而在下一轮分配中获得更多推荐。
所以 FBT 的收益不只是省下的人工，
而是\kw{把履约确定性转化为账号的长期权重}。
反过来说，选择自发货的卖家必须接受一个现实：
你必须用更高的内容质量或更低的价格，
去补偿履约不确定性带来的权重损失。
\end{insight}

\section{平台对比与模式选择矩阵}
\label{sec:plat-matrix}

\subsection{七个维度的横向对比}

把本章讨论的三个平台放在七个维度上比较，
可以避免「哪个平台更好」这种没有答案的问题，
转而问「我的能力结构匹配哪个平台」。

\begin{table}[htbp]
\centering
\caption{三大平台在七个关键维度的横向对比（定性判断，费用与规则以官方最新为准）}
\label{tab:plat-matrix}
\scriptsize
\begin{tabularx}{\textwidth}{L{1.7cm} Y Y Y}
\toprule
\tabhead{维度} & \tabhead{Temu} & \tabhead{SHEIN（自营与 Marketplace）} & \tabhead{TikTok Shop} \\
\midrule
流量获取难度 & 低（平台负责，卖家只需过核价） & 低到中（自营流量强，第三方需竞争） & 高（内容与投放能力是硬门槛） \\
\addlinespace
定价权 & 无 & 很弱（价格对齐机制） & 完整（但受竞争价格约束） \\
\addlinespace
资金占用 & 中（国内备货加账期缺口） & 中到高（备货加平台仓） & 高（海外备货加投流加达人分佣） \\
\addlinespace
毛利水平 & 最低，靠成本与周转 & 中，靠款式与供应链效率 & 最高，但退货与营销占比大 \\
\addlinespace
进入门槛 & 低（有产能即可） & 中（款式与快反能力） & 高（内容、达人、履约团队） \\
\addlinespace
合规要求 & 平台承担部分，卖家需完成产品合规 & 平台规则严格，目的国合规要求高 & 卖家承担较多，含内容与广告合规 \\
\addlinespace
政策风险 & 目的国关税与贸易政策变化 & 目的国贸易政策与平台策略调整 & 平台所有权与数据相关监管不确定性 \\
\bottomrule
\end{tabularx}
\end{table}

\subsection{多平台组合：把它当成一个投资组合问题}

很多卖家的真实问题是「几个平台同时做，资源怎么分」。
这个问题在数学上与资产组合选择同构：
每个平台贡献期望利润，也带来方差，
而平台之间的风险是\kw{正相关的}——
关税政策、消费景气、物流价格这类冲击会同时影响所有平台。
写成风险调整后的优化问题：
\begin{equation}
\max_{\mathbf{w}} \;\; \E\big[\pi(\mathbf{w})\big] - \lambda \sqrt{\Var\big[\pi(\mathbf{w})\big]}
\qquad
\text{s.t.} \;\; \sum_{j} w_{j} = 1,\;\; w_{j} \ge 0 ,
\label{eq:plat-portfolio}
\end{equation}
其中 $w_{j}$ 是投入平台 $j$ 的资源份额（可以是资金、产能或人力），
$\lambda$ 是卖家的\kw{风险厌恶系数}。
这个模型的实践价值在于给出三个反直觉结论：

\begin{enumerate}
  \item \textbf{分散化降低的是「平台特有风险」，不能降低「系统性风险」。}
        多开一个平台可以规避「某个平台突然改规则」的风险，
        但无法规避「目的国加征关税」或「海运价格翻倍」。
        真正对冲系统性风险的手段是\kw{多市场、多产地}，见 \cref{ch:supply}。
  \item \textbf{资源份额不等于成交额份额。} 用 \cref{eq:plat-portfolio} 思考时，
        要按\kw{资金占用与人力占用}分配权重，
        而不是按成交额。一个成交额占比 40\% 但占用 70\% 库存资金的平台，
        在风险模型里的权重远高于 40\%。
  \item \textbf{$\lambda$ 应该随自有资金规模下降。}
        资金充裕的卖家可以承担更高的波动去博取更高的期望利润；
        现金流紧张的卖家必须优先选择\kw{回款快、占用低}的模式，
        哪怕它的期望利润率更低。这就是为什么\kw{同一个平台，
        对不同卖家的最优性完全不同}。
\end{enumerate}

\subsection{判断路径：什么品适合什么模式}

\begin{toolbox}[模式选择的四步判断路径]
按顺序回答下面四个问题，可以在半小时内缩小到一到两个候选模式。
\begin{enumerate}
  \item \textbf{这个品有没有「内容可讲性」？}
        有强视觉冲击、使用演示、痛点场景的品（家居小工具、美妆、服饰、创意礼品）
        优先考虑内容电商；
        纯功能标品（配件、耗材、基础五金）更适合全托管或半托管。
  \item \textbf{单件毛利能不能覆盖跨境履约的固定成本？}
        用 \cref{eq:temu-cost} 与 \cref{eq:fm-profit-self} 分别测算。
        如果零售价低于「履约加关税加尾程」之和的一定倍数，
        自运营在结构上不可行，应转向全托管。
  \item \textbf{我的资金能承受多长的现金转换周期？}
        用 \cref{eq:fm-cashgap} 算出三种模式下的现金缺口。
        缺口超过可动用资金的一半时，必须先解决资金结构，再谈模式。
  \item \textbf{我愿意投入多少「非生产性人力」？}
        内容、投放、达人、客服都是非生产性岗位。
        团队里没有这类人时，强行做自运营的失败率极高；
        反过来，如果团队只有运营能力没有产能，全托管是更合理的起点。
\end{enumerate}
\textbf{组合建议}（示范性策略，非投资建议）：
用全托管承接\kw{确定性产能与基础现金流}，
用半托管承接\kw{已有海外仓能力的中高毛利款}，
用内容电商承接\kw{需要品牌化的差异化款}。
三者的资金与人力占比建议按 \cref{eq:plat-portfolio} 每季度重算一次。
\end{toolbox}

\section{平台规则与风控：违规、处罚与申诉}
\label{sec:plat-risk}

\subsection{五类高频违规与对应的处罚手段}

平台规则文本动辄数百页，但真正高频触发处罚的违规类型并不多。
把它们归纳成五类，并说清「平台在防什么」，
比背诵条款更容易形成正确的直觉。

\begin{table}[htbp]
\centering
\caption{五类高频违规、平台动因与常见处罚手段（具体条款与标准以官方最新规则为准）}
\label{tab:plat-penalty}
\small
\begin{tabularx}{\textwidth}{L{2.0cm} L{3.2cm} L{3.1cm} Y}
\toprule
\tabhead{违规类型} & \tabhead{典型表现} & \tabhead{处罚手段} & \tabhead{申诉要点} \\
\midrule
知识产权侵权 & 使用未授权商标、图案、外观或专利结构 & 下架、扣分、冻结货款，严重时封店 & 提供授权链条、设计源文件、在先使用证据 \\
\addlinespace
品控与安全 & 抽检不合格、客诉集中、存在安全隐患 & 罚款、批次退回、限制上新 & 提交第三方检测报告与批次追溯记录 \\
\addlinespace
虚假宣传 & 夸大功效、绝对化用语、图文不符 & 下架、罚款、降低权重 & 提供整改后的页面与素材举证 \\
\addlinespace
刷单与虚假交易 & 自买自评、异常订单聚集、评价操纵 & 清理评价、扣分、清零权重、冻结资金 & 提供真实物流与支付凭证、异常订单说明 \\
\addlinespace
价格与促销违规 & 虚构原价、先涨后降、承诺未兑现 & 罚款、取消活动资格 & 提供历史价格与促销备案记录 \\
\bottomrule
\end{tabularx}
\end{table}

理解这张表的关键在于：\kw{每一类违规背后都有一个平台必须保护的指标}。
侵权保护的是平台自身的法律风险，
品控保护的是退货率与消费者信任，
虚假宣传保护的是内容与搜索的可信度，
刷单保护的是数据质量（进而保护算法与广告定价），
价格违规保护的是「低价心智」这个核心资产。
当卖家理解了平台在保护什么，
就能预测规则的演进方向，而不是被动等待处罚通知。
需要提醒的是，任何被当作考核目标的指标都会承受被\kw{过度优化}的压力
（即 Goodhart 定律式的失效，见 \cite{goodhart1975}）：
平台侧的「质量分」「履约分」一旦成为唯一考核项，
卖家就会沿着代理指标与真实质量之间的缝隙优化，
这也是平台规则必须持续迭代、卖家必须持续跟进的根源。

\subsection{处罚的梯度与「可修复性」}

处罚不是二元的，它通常是一个梯度：
\kw{警告、扣分、罚款、下架、限制上新、冻结货款、封店}。
其中对卖家现金流影响最大的是\kw{冻结货款}——
它把「暂时性的规则问题」直接转化为「生存问题」。
把冻结的影响写成公式：
\begin{equation}
\text{冻结冲击}
= \frac{W_{\text{在途资金}}}{W_{\text{可动用资金}}} \times \frac{1}{\text{周转次数}} ,
\label{eq:plat-freeze}
\end{equation}
其中 $W_{\text{在途资金}}$ 是被冻结的应收款。
当周转次数较高（即生意本身依赖快速滚动）时，
同样的冻结金额造成的冲击会被放大——
\textbf{周转越快的生意，对资金冻结越没有抵抗力}。
这是很多「看起来很赚钱」的卖家突然倒下的真实原因。

\begin{toolbox}[申诉与资金解冻的实操路径]
以下是一套通用的操作框架，具体渠道、时限与材料要求\kw{以平台官方最新规则与卖家后台为准}。
\begin{enumerate}
  \item \textbf{第一时间固定证据}：导出申诉期内的订单、物流轨迹、
        质检记录、授权文件与历史价格快照。
        平台申诉的时限通常很短，过期后证据价值大幅下降。
  \item \textbf{区分「事实错误」与「规则理解差异」}：
        前者要用证据直接推翻（如物流签收记录证明非虚假发货）；
        后者要提交整改方案与承诺，而不是反复申辩。
        把两类问题混在一封申诉里，是最常见的失败模式。
  \item \textbf{走完内部渠道再考虑升级}：先在卖家后台的申诉入口提交，
        保留工单号；升级到客户经理或品类买手；
        仍无结果时再研究平台的争议解决机制。
        跳过内部渠道通常不会加快处理。
  \item \textbf{同步做现金流预案}：在等待期间就启动
        \kw{压缩备货、协商上游账期、暂停广告投放}三条动作，
        不要把希望完全寄托在解冻时间上。
  \item \textbf{把每次处罚转成流程改进项}：建立违规台账，
        记录\kw{事由、损失金额、根因、对应流程改动}，
        每季度复盘一次。同一类违规出现两次，说明流程没有真正修改。
\end{enumerate}
\end{toolbox}

\begin{pitfall}[「先用低价冲量，违规了再说」]
\textbf{现象}：为了快速起量，卖家在图片中使用知名 IP 元素、
在文案中使用绝对化功效表述、并配合短期促销制造「原价」锚点。
前两个月数据很好，第三个月集中收到侵权投诉与价格违规处罚，
账号被扣分、货款被冻结，量掉到零。\\
\textbf{根因}：这类操作的风险不是\kw{线性累积}而是\kw{阈值触发}的。
平台的风控通常基于\kw{投诉量、举报密度与事后审计}，
在小规模时几乎不会被发现，规模上去之后必然暴露；
而处罚力度与\kw{历史经营规模}相关，
所以「做得越大、罚得越重」。
卖家把前期的「没被处罚」误读为「这条路可行」，
本质上是把\kw{未被观测}当成了\kw{被允许}。\\
\textbf{排查方法}：
(1) 上新前做\kw{合规清单核对}——图片、文案、商标、专利、认证五项逐一确认，
        并留存核对人；
(2) 把「绝对化用语、功效承诺、未授权 IP 元素」做成关键词黑名单，
        接入素材审核流程；
(3) 促销价格必须保留完整的价格历史记录，避免被认定为虚构原价；
(4) 对任何「同行都在这么做」的说法保持警惕——
        同行的存活不能证明该做法合规，只能证明它暂时没被处理。
\end{pitfall}

\section{卖家生存策略：从供货商到品牌商}
\label{sec:plat-strategy}

\subsection{四条并行的生存策略}

在规则随时会变的环境里，生存策略的核心不是「押对平台」，
而是\kw{让自己在任何一种规则下都还有第二条路}。
可操作的策略有四条。

\textbf{(1) 多平台分散，但共享同一条供应链。}
分散的正确做法不是「每个平台建一套货」，
而是\kw{用同一套产品能力去适配多个渠道的准入规则}。
产品规格、包装与认证尽量做成「一次开发、多处合规」，
把渠道差异隔离在\kw{图片、文案与报价}这三件可变要素上。
这样每增加一个平台，边际成本主要是运营而不是重新开发。
反过来说，如果每上一个平台都要重新开模、重新认证，
多平台分散在经济上就是不成立的。

\textbf{(2) 差异化选品，把「同款比价」变成「无可比价」。}
核价机制的杀伤力来自「同款」这个前提。
只要你能让商品在\kw{规格、结构、材质、组合或包装}上无法被一一对应，
比价就失去了基准。可操作的差异化手段包括：
与工厂联合开发独家模具、使用独家面料或配色、
做成多件组合装（改变单位可比价）、
以及在包装与配件上形成独有特征。
这些手段的共同点是\kw{提高被复制的成本}，而不是提高售价。

\textbf{(3) 自有品牌与私域：在合规边界内沉淀资产。}
品牌化的第一步不是注册商标，而是\kw{让消费者能记住并指名购买}。
在平台体系内，这意味着稳定的视觉识别、一致的质量水平与可预期的服务。
私域导流的合规边界必须严格把握：
\kw{平台规则通常限制在站内引导用户离开平台交易}，
且各目的国对营销触达（邮件、短信、推送）有明确的法律要求，
例如前置同意机制与退订机制。
正确的做法是在\kw{平台允许的范围内}沉淀可复用的资产：
粉丝、评价、内容素材、达人关系与用户反馈数据，
并把它们沉淀在\kw{自己可长期访问}的载体上。
任何试图绕开平台规则把用户导走的做法，
都会在某个时点触发封店，而封店会同时清空你所有的平台侧资产。

\textbf{(4) 从供货商升级为品牌商：一条有顺序的路径。}
升级不是「注册个商标然后涨价」，而是能力的逐层叠加。

\begin{figure}[htbp]
\centering
\begin{tikzpicture}[
  node distance=6mm,
  st/.style={draw=NavyMid, fill=BoxBG, rounded corners=2pt, minimum height=9mm,
             inner sep=4pt, text width=2.9cm, align=center, font=\sffamily\tiny},
  capnode/.style={draw=TealDeep, fill=BoxBGAlt, rounded corners=1pt, inner sep=3pt,
              text width=2.9cm, align=left, font=\sffamily\tiny},
  note/.style={draw=RuleGray, fill=white, rounded corners=2pt, inner sep=4pt,
               text width=13.4cm, align=left, font=\sffamily\tiny}
]
\node[st] (a) at (0,0) {\textbf{供货商}\\按图生产、按单交付};
\node[st] (b) at (3.6,0) {\textbf{白牌商家}\\自主定价、自主履约};
\node[st] (c) at (7.2,0) {\textbf{自有品牌}\\设计、命名、复购};
\node[st] (d) at (10.8,0) {\textbf{品牌商}\\多市场、全渠道};
\node[capnode] at (0,-1.35) {能力：成本、交期、品控};
\node[capnode] at (3.6,-1.35) {能力：选品、投放、内容};
\node[capnode] at (7.2,-1.35) {能力：产品定义、用户运营};
\node[capnode] at (10.8,-1.35) {能力：渠道管理、合规、组织};
\node[note] at (5.4,-2.9)
  {每一步都不是「加一个头衔」，而是\emph{新增一种能力并承担一类新风险}。
   跳过中间步骤（例如从纯供货直接做品牌）通常失败在
   「既没有需求获取能力，也没有承担库存与合规风险的现金流」上。};
\draw[-{Stealth}, thick, SlateGray] (a) -- (b);
\draw[-{Stealth}, thick, SlateGray] (b) -- (c);
\draw[-{Stealth}, thick, SlateGray] (c) -- (d);
\end{tikzpicture}
\caption{卖家从供货商升级为品牌商的能力叠加路径}
\label{fig:plat-upgrade}
\end{figure}

\begin{insight}[「品牌」在平台体系里的现实定义]
在很多卖家的理解里，品牌等于注册商标加好看的主图。
但平台体系里的品牌其实只有三个可观测的表现：
\kw{(1) 有搜索量}（用户主动搜你的名字），
\kw{(2) 有复购}（同一用户二次购买，且不依赖新的广告触达），
\kw{(3) 有溢价能力}（同规格商品能比白牌贵，且销量不掉）。
这三个指标都可以在后台数据里被验证，
也都无法通过「重新做一套视觉」伪造。
用它们来定义品牌目标，可以让品牌建设从一句口号
变成一组可以按季度考核的指标。
\end{insight}

\subsection{用数据把多平台经营管起来}

多平台经营最大的风险不是竞争，而是\kw{算不清账}。
当同一个 SKU 同时供给三个渠道时，
很容易出现「总收入增长、净利润下降」的局面，
原因是高毛利渠道的量被低毛利渠道的备货挤占了。
解决办法是建立\kw{按「渠道乘 SKU」粒度的单位经济学看板}，
把 \cref{eq:fm-profit,eq:fm-profit-self,eq:temu-penalty} 中的项目全部落到字段上。
下面这段查询演示了最小可用的核算结构（示意，字段名需按实际系统调整）。

\begin{lstlisting}[style=growth, language=SQLx, caption={按渠道与 SKU 核算单件净利（示意）}]
SELECT
    channel,
    sku_id,
    SUM(units)                       AS units,
    SUM(revenue)                     AS revenue,
    SUM(cogs)                        AS cogs,
    SUM(logistics + tariff + warehouse)
                                     AS fulfilment_cost,
    SUM(ad_spend + creator_commission)
                                     AS traffic_cost,
    SUM(penalty + refund_loss)       AS rule_cost,
    SUM(revenue - cogs - logistics - tariff - warehouse
        - ad_spend - creator_commission
        - penalty - refund_loss)     AS gross_profit,
    SUM(revenue - cogs - logistics - tariff - warehouse
        - ad_spend - creator_commission
        - penalty - refund_loss) / NULLIF(SUM(units), 0)
                                     AS unit_profit,
    SUM(inventory_days * cogs) / NULLIF(SUM(cogs), 0)
                                     AS avg_inventory_days
FROM fact_order_sku_daily
WHERE stat_date >= DATE_SUB(CURRENT_DATE, INTERVAL 90 DAY)
GROUP BY channel, sku_id
HAVING SUM(units) > 0
ORDER BY unit_profit ASC;
\end{lstlisting}

这段查询的输出按 \code{unit\_profit} 升序排列，
也就是\kw{把最亏钱的渠道与 SKU 排在最前面}——
这是最有价值的排序方式，因为经营改善的边际收益在亏损项上最高。
配合 \cref{eq:shein-turn} 的周转指标，
就能同时看到「赚不赚钱」与「钱转得快不快」这两件事。

\section{本章小结与下一步}

\begin{keypoint}[本章要点回顾]
\begin{enumerate}
  \item 三种平台模式的本质区别是\kw{决策权与风险的归属}：
        全托管卖产能（供货价结算、平台定价与履约），
        半托管卖海外库存运营能力，自主运营卖需求获取能力。
        选型的起点永远是「我已经具备什么能力」，而不是「哪个模式毛利高」。
  \item 全托管单件利润由毛差、赔付率与账期共同决定（\cref{eq:fm-profit,eq:fm-net}）；
        账期对现金流的影响可以用现金转换周期量化（\cref{eq:fm-cashgap}）——
        \textbf{利润变成应收款，是很多工厂「看起来赚钱但活不下去」的真实原因}。
  \item 核价是「倒推」而不是「议价」（\cref{eq:temu-quote}），
        且通过「同类最低价」形成\kw{单向棘轮}，长期毛利预期必须按逐步下行建模。
        卖家报价的真实下限要由 \cref{eq:temu-cost} 计算，
        所有按金额比例发生的损失必须放在分母里。
  \item 罚款由按单、按批、按金额比例三种形式组合（\cref{eq:temu-penalty}）：
        罚款率与报价同向放大、批次数越多总罚款越高、
        固定罚款对小单价品类的伤害是致命的。
  \item JIT 不消灭库存风险，它把「成品库存风险」换成了「产能预留成本」；
        是否更划算取决于两者在同一量纲下的比较。
  \item 小单快反的核心是\kw{用真实销售数据替代需求预测}（\cref{eq:shein-batch}），
        其临界售罄率（\cref{eq:shein-threshold}）说明：
        单件毛差越小、首单量越小，测款的成功门槛越高，
        因此低毛利标品在结构上不适合照搬这套打法。
  \item 库存周转是财务指标：\kw{资金回报率等于单件毛利率乘以年周转次数}
        （\cref{eq:shein-turn}）。
        低毛利业务只有两条出路——转得更快，或每转一圈赚得更多。
  \item 内容电商的成交漏斗可以完整分解为
        \cref{eq:tt-live-gmv} 的六个因子，
        「人货场」正好对应其中的进入率、点击与转化、客单价与退款率。
  \item 达人分销的盈亏平衡佣金率由 \cref{eq:tt-breakeven} 给出；
        多数「佣金谈不拢」的问题，根因是退款率与履约成本被低估。
        FBT 的价值不只是省事，而是把履约确定性转化为账号的长期权重。
  \item 平台选择是一个组合优化问题（\cref{eq:plat-portfolio}）：
        分散化能降低平台特有风险，但不能降低关税与物流这类系统性风险；
        资源份额应按资金与人力占用分配，而不是按成交额分配。
  \item 五类高频违规背后各自对应平台必须保护的指标；
        处罚是梯度而非二元，\kw{资金冻结}是最致命的形态
        （\cref{eq:plat-freeze}），且周转越快的生意越扛不住冻结。
  \item 生存策略是四条并行：多平台共享同一条供应链、
        用差异化让比价失去基准、在合规边界内沉淀品牌与数据资产、
        按「供货商、白牌商家、自有品牌、品牌商」逐层升级能力。
        品牌的可验证定义只有三个：有搜索、有复购、有溢价。
\end{enumerate}
\end{keypoint}

\begin{interview}[本章对应的面试题]
\begin{enumerate}
  \item 全托管与半托管的本质区别是什么？
        如果平台的定价权不变，为什么把物流责任转移给卖家会改变卖家的能力要求？
  \item 一个供货价 18 元、完全成本 12 元的商品，
        在赔付率 3\%、退货率 6\%、账期 45 天、年化资金成本 8\% 的假设下，
        单件净利是多少？请写出计算过程并指出哪个假设最敏感。
  \item 平台的核价机制为什么是「单向棘轮」？
        这种机制对品类的长期毛利结构与卖家数量会有什么影响？
  \item JIT 与备货各自的成本结构是什么？
        在什么条件下 JIT 反而比备货更贵？
  \item SHEIN 的小单快反为什么难复制？
        请从产业带密度、账期结构、买手组织与临界售罄率四个角度回答。
  \item 达人分销中，如何计算可承受的最高佣金率？
        如果退款率从 8\% 上升到 15\%，可承受佣金率会如何变化，
        你会用什么手段来对冲？
  \item 多平台经营时，资源应该按成交额份额分配还是按资金占用分配？为什么？
        什么样的风险是分散化无法对冲的？
  \item 平台的处罚梯度中，为什么「冻结货款」比「罚款」更危险？
        请给出一个针对资金冻结的 30 天应急预案。
\end{enumerate}
参考答案与答题框架见 \cref{app:interview}。
\end{interview}

到这里，平台侧的游戏规则已经讲清楚：
它决定了谁能卖、以什么价格卖、钱什么时候到手、出问题谁赔。
但平台模式只是\kw{渠道结构}，
真正决定一个卖家能不能长期赚钱的，是渠道背后的\kw{商品与供应链}——
选什么品、找哪个供应商、怎么定安全库存、
头程与海外仓怎么选、退货怎么处理、牛鞭效应怎么抑制。
这些问题与平台规则无关，却决定了同一套规则下不同卖家的天壤之别。
下一章 \cref{ch:supply} 我们就进入这一层：
选品打分模型、供应商评估与谈判、品控与认证、
安全库存与补货规则、履约与逆向物流，
以及关税与产地多元化冲击下的供应链韧性设计。
读完那一章再回看本章的 \cref{eq:temu-cost,eq:shein-threshold}，
你会发现它们的分母与分子，最终都由供应链能力决定。
